\section{Continuous Microlocal Derivation}
\label{app:continuous}

This appendix records the continuous microlocal layer in more detail than the main text. The main body uses a dual spine: a short-time half-wave off-graph decay theorem and an abstract canonical-graph packet sparsity theorem. This appendix gives the lemma chain behind those statements.

For the single flagship manuscript, the FoCM-style proof map is:
\[
\begin{gathered}
\text{short-time parametrix}
\Rightarrow
\text{restricted non-stationary phase} \\
\Rightarrow
\text{packet almost-diagonalization}
\Rightarrow
\text{canonical-tube sparsity} \\
\Rightarrow
\text{transported-symbol reduction}
\Rightarrow
\text{network-realizable \MiNO-Core theorem}.
\end{gathered}
\]
The finite landing layers of this map are machine-checked. The continuous arrows are classical manuscript mathematics restricted to Euclidean compact phase windows, Gaussian/Gabor packet frames, finite branch count, and a no-caustic short-time half-wave regime. This division gives continuous microlocal structure for the architecture and finite formalization for the computational landing layer.

\subsection{Short-time half-wave parametrix}

Let
\[
p(x,\xi)=c(x)|\xi|,
\qquad
0<c_{\min}\le c(x)\le c_{\max}<\infty,
\]
with $c\in C^\infty(\R^d)$. On a short time interval $|t|\le T$ chosen before caustic formation on the frequency band seen by the packet frame, the positive half-wave propagator $U_+(t)$ admits a standard oscillatory representation
\[
U_+(t)u(x)
\sim
(2\pi h)^{-d}
\iint e^{\frac{i}{h}(\Phi_t(x,\xi)-y\cdot\xi)}
a_t(x,\xi;h)\,u(y)\,dy\,d\xi,
\]
where:
\begin{enumerate}
\item $\Phi_t$ generates the canonical graph of the Hamiltonian flow $\kappa_t$ associated with $p$;
\item $a_t$ is a classical amplitude with uniform symbol bounds on compact time intervals;
\item the graph remains single-valued on the chosen time window.
\end{enumerate}
This is the exact place where we restrict to short time and Euclidean local coordinates. The theorem does not try to address global caustic geometry.

\subsection{Packet testing and localized phase}

Let $\varphi_\eta^h$ and $\varphi_\gamma^h$ be analytic Gaussian/Gabor-type packets centered at $z_\eta=(x_\eta,\xi_\eta)$ and $z_\gamma=(x_\gamma,\xi_\gamma)$. Inserting the packets into the parametrix gives
\[
M_{\gamma\eta}^h(t)
:=
\langle U_+(t)\varphi_\eta^h,\varphi_\gamma^h\rangle
=
(2\pi h)^{-d}\iiint
e^{\frac{i}{h}\Psi_{t,\gamma,\eta}(x,y,\xi)}
b_{t,\gamma,\eta}(x,y,\xi;h)\,dx\,dy\,d\xi,
\]
where the localized phase $\Psi_{t,\gamma,\eta}$ is the half-wave phase plus the packet phases, and $b_{t,\gamma,\eta}$ is a compactly localized amplitude coming from the packet windows and the FIO amplitude.

The continuous packet problem is therefore reduced to an oscillatory-integral question: how does the integral behave when $z_\gamma$ is far from $\kappa_t(z_\eta)$?

\subsection{Phase-gradient lower bound outside the canonical tube}

The key microlocal lemma is a non-stationarity statement.

\begin{lemma}[Phase-gradient lower bound]
\label{lem:phase-gradient}
Fix a tube radius parameter $\rho>0$. There exist constants $c_\rho,C_\rho>0$ such that if
\[
\mathrm{dist}\bigl(z_\gamma,\kappa_t(z_\eta)\bigr)\ge \rho h^{1/2},
\]
then on the support of the localized amplitude one has
\[
|\nabla_{x,y,\xi}\Psi_{t,\gamma,\eta}(x,y,\xi)|
\ge
c_\rho\,\mathrm{dist}\bigl(z_\gamma,\kappa_t(z_\eta)\bigr)
- C_\rho h^{1/2}.
\]
In particular, after enlarging $\rho$ if necessary,
\[
|\nabla_{x,y,\xi}\Psi_{t,\gamma,\eta}(x,y,\xi)|
\gtrsim
\mathrm{dist}\bigl(z_\gamma,\kappa_t(z_\eta)\bigr)
\]
uniformly outside the canonical tube.
\end{lemma}

\begin{proof}[Proof idea]
The packets localize the variables to neighborhoods of size $h^{1/2}$ around $(x_\gamma,\xi_\gamma)$ and $(x_\eta,\xi_\eta)$. On such neighborhoods, the stationary point equations for $\Psi_{t,\gamma,\eta}$ reduce to the canonical relation for the half-wave phase. If $z_\gamma$ is outside the tube around $\kappa_t(z_\eta)$, no stationary point survives on the localized support. Smooth dependence of the parametrix on $(x,\xi)$ and the no-caustic hypothesis then yield a quantitative lower bound.
\end{proof}

Lemma~\ref{lem:phase-gradient} is the PDE-specific step that replaces the old abstract assumption. Once it is available, the remainder of the packet argument is a classical non-stationary phase calculation.

\subsection{Restricted non-stationary phase bridge}

For the packet theorem used here, the full continuous calculus is not needed at the point where packet almost-diagonalization begins. The exact bridge can be isolated as follows.

\begin{proposition}[Restricted packet non-stationary phase theorem]
\label{prop:restricted-nsp-appendix}
Consider a packet-localized oscillatory integral representation
\[
M_{\gamma\eta}^h
=
\int e^{\frac{i}{h}\Psi_{\gamma\eta}^h(z)} b_{\gamma\eta}^h(z)\,dz
\]
on a compact phase-space window.  After the packet rescaling
$z=z_{\gamma\eta}^0+h^{1/2}\zeta$, write the integral as
\[
M_{\gamma\eta}^h
=
\int_{\R^m} e^{i\Theta_{\gamma\eta}^h(\zeta)}
  a_{\gamma\eta}^h(\zeta)\,d\zeta,
\]
where $a_{\gamma\eta}^h$ is compactly supported in a fixed compact set
independent of $h,\gamma,\eta$.  Set
\[
D_{\gamma\eta}
:=
1+h^{-1/2}\mathrm{dist}(z_\gamma,\kappa(z_\eta)).
\]
Fix $N\ge 1$.  Assume:
\begin{enumerate}
\item the phase comes from a single canonical graph on the chosen short-time window;
\item the rescaled phase has no stationary point off the canonical tube and satisfies, for all multiindices $1\le |\alpha|\le N+1$,
\[
\bigl|\partial_\zeta^\alpha \Theta_{\gamma\eta}^h(\zeta)\bigr|
\le C_\alpha D_{\gamma\eta},
\qquad
\bigl|\nabla_\zeta \Theta_{\gamma\eta}^h(\zeta)\bigr|
\ge cD_{\gamma\eta};
\]
\item the rescaled localized amplitude satisfies the uniform $L^1$ seminorm bound
\[
\sum_{|\alpha|\le N}
\|\partial_\zeta^\alpha a_{\gamma\eta}^h\|_{L^1}
\le B_N .
\]
\end{enumerate}
Then there is a constant $C_N'$ depending only on $N$, $c$, the phase derivative constants $C_\alpha$, the dimension, and the fixed support set such that
\[
|M_{\gamma\eta}^h|
\le
C_N' B_N
\bigl(1+h^{-1/2}\mathrm{dist}(z_\gamma,\kappa(z_\eta))\bigr)^{-N}.
\]
\end{proposition}

\begin{proof}
Let
\[
v(\zeta)
:=
\frac{\nabla_\zeta\Theta_{\gamma\eta}^h(\zeta)}
     {|\nabla_\zeta\Theta_{\gamma\eta}^h(\zeta)|^2},
\qquad
L:=\frac{1}{i}v(\zeta)\cdot\nabla_\zeta .
\]
Then $Le^{i\Theta_{\gamma\eta}^h}=e^{i\Theta_{\gamma\eta}^h}$.  Since the amplitude is compactly supported in a fixed compact set, repeated integration by parts gives
\[
M_{\gamma\eta}^h
=
\int e^{i\Theta_{\gamma\eta}^h(\zeta)}
(L^\ast)^N a_{\gamma\eta}^h(\zeta)\,d\zeta .
\]

It remains to bound $(L^\ast)^N a_{\gamma\eta}^h$ in $L^1$.  The assumptions imply a uniform coefficient estimate for the vector field $v$: for every multiindex $|\alpha|\le N$,
\[
|\partial_\zeta^\alpha v(\zeta)|
\le
C_{\alpha,c}\,D_{\gamma\eta}^{-1}.
\]
Indeed, $v$ is a rational expression in the first derivatives of
$\Theta_{\gamma\eta}^h$ with denominator
$|\nabla\Theta_{\gamma\eta}^h|^2$; differentiating $v$ produces sums of terms whose numerators contain derivatives of $\Theta_{\gamma\eta}^h$ bounded by $C_\alpha D_{\gamma\eta}$ and whose denominators contain at least one net power of
$|\nabla\Theta_{\gamma\eta}^h|$.  The lower bound
$|\nabla\Theta_{\gamma\eta}^h|\ge cD_{\gamma\eta}$ therefore leaves one factor $D_{\gamma\eta}^{-1}$ in every coefficient derivative used below.

The formal adjoint has the form
\[
L^\ast f
=
i^{-1}\bigl[-v\cdot\nabla f-(\div v)f\bigr],
\]
up to the harmless sign convention for complex conjugation.  Thus $L^\ast$ is a first-order differential operator whose coefficients and coefficient derivatives up to order $N-1$ are $O(D_{\gamma\eta}^{-1})$.  Expanding $(L^\ast)^N$ by induction yields
\[
(L^\ast)^N a_{\gamma\eta}^h
=
\sum_{|\alpha|\le N} c_{\alpha,N,\gamma,\eta}(\zeta)
\partial_\zeta^\alpha a_{\gamma\eta}^h(\zeta),
\qquad
\|c_{\alpha,N,\gamma,\eta}\|_{L^\infty}
\le
C_N''D_{\gamma\eta}^{-N}.
\]
Consequently,
\[
\|(L^\ast)^N a_{\gamma\eta}^h\|_{L^1}
\le
C_N''D_{\gamma\eta}^{-N}
\sum_{|\alpha|\le N}\|\partial_\zeta^\alpha a_{\gamma\eta}^h\|_{L^1}
\le
C_N''B_ND_{\gamma\eta}^{-N}.
\]
Taking absolute values in the integrated-by-parts representation proves the stated estimate.
\end{proof}

\subsection{Repeated integration by parts}

Define the standard non-stationary phase operator
\[
L
:=
\frac{\nabla\Psi_{t,\gamma,\eta}\cdot\nabla_{x,y,\xi}}
{i|\nabla\Psi_{t,\gamma,\eta}|^2}.
\]
By construction,
\[
L e^{\frac{i}{h}\Psi_{t,\gamma,\eta}}
=
h^{-1} e^{\frac{i}{h}\Psi_{t,\gamma,\eta}}.
\]
Integrating by parts $N$ times gives
\[
M_{\gamma\eta}^h(t)
=
h^N
(2\pi h)^{-d}
\iiint
e^{\frac{i}{h}\Psi_{t,\gamma,\eta}}
(L^\ast)^N b_{t,\gamma,\eta}\,dx\,dy\,d\xi.
\]
Lemma~\ref{lem:phase-gradient} shows that every application of $L^\ast$ contributes one inverse power of the tube distance, while analyticity and bounded overlap of the packet windows keep differentiated amplitudes summable. This yields
\[
|M_{\gamma\eta}^h(t)|
\le
C_{T,N}
\Bigl(1+h^{-1/2}\,\mathrm{dist}\bigl(z_\gamma,\kappa_t(z_\eta)\bigr)\Bigr)^{-N},
\]
which is Theorem~\ref{thm:wave-offgraph}.

\subsection{Shell decomposition for abstract FIO sparsity}

Fix a source packet $\eta$ and write
\[
\mathcal S_\ell(\eta)
=
\Bigl\{
\gamma\in\Gamma_h :
2^\ell h^{1/2}
\le
\mathrm{dist}\bigl(z_\gamma,\kappa(z_\eta)\bigr)
<
2^{\ell+1} h^{1/2}
\Bigr\},
\qquad \ell\ge 0.
\]
Assume the canonical-graph packet matrix satisfies:
\begin{enumerate}
\item \emph{off-graph decay:}
\[
|F_{\gamma\eta}^h|
\le
C_N 2^{-N\ell}
\qquad\text{for }\gamma\in\mathcal S_\ell(\eta);
\]
\item \emph{shell counting:}
\[
\#\mathcal S_\ell(\eta)\le C_{\mathrm{cnt}}2^{q\ell}.
\]
\end{enumerate}
Then the shell mass satisfies
\[
\sum_{\gamma\in\mathcal S_\ell(\eta)} |F_{\gamma\eta}^h|
\le
C 2^{(q-N)\ell}.
\]
Summing shells $\ell\ge L$ gives
\[
\sum_{\ell\ge L}\sum_{\gamma\in\mathcal S_\ell(\eta)} |F_{\gamma\eta}^h|
\le
C' 2^{(q-N)L}.
\]
The number of packets retained through shell $L-1$ is $O(2^{qL})$, so if $K$ is the retained stencil size then $K\asymp 2^{qL}$ and therefore
\[
2^{(q-N)L}\asymp K^{1-N/q}.
\]
This is the entire proof of Theorem~\ref{thm:packet-sparsity}.

\subsection{Restricted packet almost-diagonalization bridge}

Once Proposition~\ref{prop:restricted-nsp-appendix} and the shell-counting hypothesis are available, the remaining step toward a theorem-grade packet law is local. Near the canonical graph, stationary phase identifies the retained entries of the packet matrix with sampled amplitudes; away from the graph, the dyadic shell tail is summable. The combination gives a restricted packet almost-diagonalization theorem:
\begin{enumerate}
\item the far-graph packet mass contributes a $K^{1-N/q}$ tail;
\item the retained tube contributes a local sampled-symbol law plus frame, asymptotic, and covering-defect remainders;
\item together they imply the bounded-stencil transported-symbol decomposition of Theorem~\ref{thm:continuous-reduction}.
\end{enumerate}
The key point is structural. The almost-diagonalization statement required by MiNO is not a vague appeal to the literature. It is the explicit sum of one restricted non-stationary-phase estimate and one local stationary-phase estimate.

\subsection{Bounded-stencil transported-symbol reduction}

The shell argument only tells us that far-graph entries are small. To get the actual \MiNO\ reference law one still needs a near-graph approximation on the retained tube. This is a stationary-phase calculation rather than a non-stationary one.

Let $\Pi_{h,K}^\star(t,\gamma)$ denote the $K$ source packets whose propagated centers lie nearest to $z_\gamma$. For $\eta\in\Pi_{h,K}^\star(t,\gamma)$, stationary phase around the localized canonical point gives
\[
\langle U_+(t)\varphi_\eta^h,\varphi_\gamma^h\rangle
=
s_{h,\gamma\eta}^\star(t)
+ O(h^{1/2}(1+\delta_h)),
\]
where $s_{h,\gamma\eta}^\star(t)$ is the sampled FIO amplitude. Hence
\[
K_h(t)
=
{\transportref_h}^{(K)}(t)+R_h^{(K,M)}(t),
\qquad
({\transportref_h}^{(K)}(t)a)_\gamma
=
\sum_{\eta\in\Pi_{h,K}^\star(t,\gamma)} s_{h,\gamma\eta}^\star(t)a_\eta,
\]
with residual consisting of:
\begin{enumerate}
\item the discarded far-shell tail from Theorem~\ref{thm:packet-sparsity};
\item the local stationary-phase remainder $O(h^M)$;
\item the packet covering mismatch $O(\delta_h)$.
\end{enumerate}
This is exactly Theorem~\ref{thm:continuous-reduction}.

\subsection{Identity-canonical pseudodifferential branch}

The elliptic side of the restricted calculus uses the same packet argument with the canonical graph replaced by the identity relation.  Let $A_h=\Op_h(a)$ with $a\in S^m_{1,0}$ on the compact phase window.  Standard pseudodifferential packet almost-diagonalization gives, for every $N\ge 1$,
\[
|\langle A_h\varphi_\eta^h,\varphi_\gamma^h\rangle|
\le
C_N \langle \xi_\eta\rangle^m
\Bigl(1+h^{-1/2}\mathrm{dist}(z_\gamma,z_\eta)\Bigr)^{-N}.
\]
Thus the relevant tube is not a transported tube around $\kappa(z_\eta)$ but the identity tube around $z_\eta$.  The same row/column shell-counting proof and Schur transfer then yield
\[
\|W_hA_hW_h^\ast-(W_hA_hW_h^\ast)^{(K_{\mathrm{id}})}\|_{\ell^2\to\ell^2}
\le
C_{\Psi}K_{\mathrm{id}}^{1-N/q},
\]
after absorbing the symbol-order envelope into the compact-window constant or into a weighted packet norm.  This is the mathematical reason the elliptic branch of \MiNO\ should be local-symbol dominated rather than transport dominated.

\subsection{Dissipative pseudodifferential branch}

For parabolic families, the microlocal law is not merely identity-canonical; it is also dissipative.  In the model case where the generator has principal symbol $q(x,\xi)$ with $\operatorname{Re}q(x,\xi)\ge c|\xi|^m$ on the high-frequency window, the short-time semigroup has principal packet action
\[
m_t(x,\xi)\approx \exp[-tq(x,\xi)].
\]
The executable branch therefore parameterizes a multiplier of the form
\[
m_\theta(x,\xi,\Delta t)
=
\exp[-\Delta t\,\operatorname{softplus}(q_\theta(x,\xi))],
\]
which enforces $0<m_\theta\le 1$ before any learned residual correction.  This is not a global Fourier fallback.  It is the dissipative-symbol regime of the same H\"ormander-style packet calculus.

\subsection{Unified branch summation}

The restricted calculus theorem in Section~\ref{sec:unified-calculus} is obtained by summing three packet laws:
\begin{enumerate}
\item FIO branches: off-canonical-tube decay plus shell counting;
\item $\Psi$DO branch: off-identity-tube decay plus shell counting;
\item smoothing branch: compact-window residual controlled by smoothing order.
\end{enumerate}
The finite counterpart is simple: once each branch has an $\ell^1$ budget, the total packet error is bounded by their sum. This summation step is machine-checked in \codepath{UnifiedCalculus.lean}; the continuous FIO and $\Psi$DO almost-diagonalization inputs are classical.

\subsection{Lean-friendly lemma chain}

The oscillatory integral proof is classical. The decomposition below gives each continuous layer a finite mirror.

\paragraph{Layer 1: canonical shells and tubes.}
\codepath{CanonicalTube.lean} defines the finite shell and tube objects and proves the basic inclusion lemmas used by any truncation argument.

\paragraph{Layer 2: shell counting interface.}
\codepath{ShellCount.lean} isolates shell-cardinality bounds from the continuous geometry that produced them. This is the place where Euclidean phase-space counting enters the finite theorem chain.

\paragraph{Layer 3: geometric shell tails.}
\codepath{PacketTail.lean} proves the finite geometric tail estimate from shell mass bounds.

\paragraph{Layer 4: abstract sparsity bridge.}
\codepath{AbstractSparsity.lean} packages shell counting plus entry decay into a canonical-tube truncation theorem.

\paragraph{Layer 5: restricted bridge wrapper.}
\codepath{RestrictedBridge.lean} packages the restricted off-graph tail with the finite transfer wrapper. This is the Lean-facing mirror of the restricted non-stationary-phase / almost-diagonalization bridge used in the manuscript.

\paragraph{Layer 6: wave-to-transfer wrapper.}
\codepath{WaveTransfer.lean} turns the truncation error into the same language as the explicit cross-resolution theorem.

The wave/FIO analysis is classical. Once it is reduced to shell counts and off-graph decay, the theorem prover handles the finite layer.

\subsection{Classical continuous inputs}

The classical inputs are:
\begin{enumerate}
\item The short-time half-wave parametrix and the identification of the canonical graph are classical.
\item The phase-gradient lower bound of Lemma~\ref{lem:phase-gradient} is classical.
\item The restricted repeated-integration-by-parts argument is proved at manuscript level in Proposition~\ref{prop:restricted-nsp-appendix}.
\item The near-graph stationary-phase identification of the retained amplitude samples is classical.
\end{enumerate}
These inputs feed the finite \MiNO\ theorem spine through explicit shell/tube estimates.

\subsection{Scope of the continuous theorem}

The theorem chain has the following scope.

\paragraph{Short time only.}
The half-wave theorem is stated only before caustics and branch crossing. Long-time propagation will generally require multibranch packet laws and more delicate geometry.

\paragraph{Euclidean local charts.}
The current proof uses $\R^d$ packet frames. Manifold generality would require additional chart and partition-of-unity bookkeeping.

\paragraph{Canonical graph regime.}
The abstract sparsity theorem assumes a single canonical graph. Multigraph or reflected propagators can still fit the framework, but only after splitting into branches.

\paragraph{Bounded-$K$ rather than dense transport.}
The whole point of the dual spine is to show that bounded stencils are natural. If the true packet matrix is structurally dense, then the transported-symbol thesis is simply the wrong primitive for that operator family.

\subsection{FoCM-strength extensions}

The extension path is:
\begin{enumerate}
\item enlarge the finite theorem core from $K=1$ to general bounded-$K$ stencils;
\item formalize the shell-counting law for a concrete Euclidean packet frame;
\item replace the restricted bridge used here by a genuinely continuous Lean proof of Proposition~\ref{prop:restricted-nsp-appendix};
\item extend the half-wave theorem to multibranch or reflected settings.
\end{enumerate}
These steps target the remaining gap between the restricted manuscript calculus and a fully formal continuous development.
